\chapter{自我与人格}

\section{自我域的职责与表示}

自我（Self）是完整主体档案中的高级认知领域，描述主体如何长期理解自身身份、角色、能力边界、偏好、价值倾向和经历。当前目标包含数种表示，不要求分别建立多个顶层子系统：

\begin{description}
\item[自我模型（Self Model）] 关于角色、能力和边界的可修订陈述。
\item[倾向与人格风格（Dispositions / Persona）] 相对稳定的特质、偏好、价值倾向、互动风格和习惯性适应。
\item[叙事身份（Narrative Identity）] 主体对自身变化、连续性与未来方向的长期叙述。
\end{description}

人格研究区分特质、特征性适应和生命叙事；自传记忆研究讨论自传知识与当前目标、自我工作模型之间的关系。这些研究帮助区分问题结构，不直接决定 Nous Wave 的数据表或学习算法。\cite{mcadams2006,mcadamsmclean2013,conwaypleydellpearce2000}

\section{认知种子与系统宪制}

初始化主体时，宿主可提供认知种子（Cognitive Seed），用于建立初始自我认知、倾向或叙事主张。系统宪制（System Constitution）保存产品约束和不能由普通认知绕过的规范规则。两类内容分别建模和授权。\cite{nous-research-package}

\StatusTag{待定}{认知种子中可演化内容与构成性约束的划分，以及覆盖或重释的管理权限。}

本节按研究资料包界定表示范围。模型和服务提供方身份归入运行配置；用户提出“以后讨厌某事”时，先记录来源与观察事件，再由自我领域依据提案和当前版本决定是否形成修订。

\section{当前实现状态}

Nous Wave 源树包含 Character Seed 和记忆、运行时基础实现；自我和人格的权威领域列入后续研究。Heptalogos 当前主体范围聚焦交互与行为管理，高级主体认知由后续跨库集成承接。
